\documentclass[11pt]{article}
\usepackage[T1]{fontenc}
\usepackage{amsmath,amssymb,amsthm}
\usepackage[margin=1in]{geometry}
\usepackage[hidelinks]{hyperref}
\newtheorem{theorem}{Theorem}
\newtheorem{lemma}[theorem]{Lemma}
\newtheorem{proposition}[theorem]{Proposition}
\newtheorem{corollary}[theorem]{Corollary}
\newcommand{\m}{\mathfrak m}
\newcommand{\C}{\mathcal C}
\newcommand{\D}{\mathcal D}
\newcommand{\ord}{\operatorname{ord}}
\newcommand{\Spec}{\operatorname{Spec}}
\setlength{\parindent}{0pt}
\setlength{\parskip}{0.5em}
\emergencystretch=3em
\title{A continuous obstacle with no greatest geometric envelope\\
Normal crossing germs, conductor fibres, and an exact canonical filtration}
\author{Independent ROOT working paper; ORIGINAL P12.3 remains open}
\date{9 October 2026}
\begin{document}
\maketitle

\begin{abstract}
Over any field and in every dimension at least two, we construct a continuous
signed obstacle on the normalized semivaluation boundary of a normal crossing
germ whose geometric continuous and decreasing classes have no greatest
minorant. The raw supremum is explicitly discontinuous and is not upper
semicontinuous. Its upper semicontinuous regularization is not geometrically
positive. We compute every canonical ideal and every canonical profile; their
pointwise convergence has an exact nonvanishing uniform error. On the
normalization, the same obstacle has a continuous geometric envelope, which
fails constancy on the infinite-conductor fibres. The only incoming theorem
is the entire qualified geometric comparison grant G331. This is a complete
scoped negative result, subject to independent review, rather than closure
of the general normal-germ part of ORIGINAL P12.3.
\end{abstract}

\section{Precise geometric input and statement}

We use the entire qualified G331 construction and comparison theorem, with
all its source, model, support, and historical-reading qualifications.
For an excellent reduced local ring and a proper primary gauge, its
independently geometric core consists of evaluations of actual negative
effective finite real Cartier divisors on actual primary blowups, with
nonnegative degree on \emph{every} contracted proper integral curve.
Its finite continuous uniform completion is \(\C_R\), and its actual
decreasing-net completion, with bottom included, is \(\D_R\).
The comparison in G331, not a new definition of positivity here, gives
\[
 \C_R=\overline{Q_R}^{\|\cdot\|_\infty},\qquad
 Q_R=\left\{-\sum_{j=1}^{r}\lambda_jv(I_j):
              \lambda_j\geq0,\ I_j\text{ unit or proper primary}\right\},
 \qquad \D_R=\{\inf_i g_i:g_i\in\C_R\text{ decrease}\}.
\]
Raw geometric and ideal cores need not be equal. We use only their proved
equality after uniform and decreasing completion. In particular, a numerical
inequality valid for every member of \(Q_R\) passes to \(\C_R\) and then
to \(\D_R\). Every member of \(\D_R\) is upper semicontinuous, being an
actual infimum of continuous functions. No increasing-limit closure is used.

Fix any field \(k\), an integer \(d\geq2\), and real numbers \(0<a<b\).
Put \(\mathbf z=(z_1,\ldots,z_{d-1})\) and
\[
 R=k[x,y,\mathbf z]_{(x,y,\mathbf z)}/(xy),\qquad
 \m=(x,y,\mathbf z),\qquad X=V_{\m}(R).
\]
Every prime support is retained. Coordinates may have value \(+\infty\).
Set
\[
 X_1=\{v:v(y)=+\infty\},\quad X_2=\{v:v(x)=+\infty\},
 \quad Z=X_1\cap X_2.
\]
Define, with \(1/(+\infty)=0\),
\[
 f(v)=\begin{cases}
       -a,&v\in X_1,\\
       -a-(b-a)/v(y),&v\in X_2.
       \end{cases}
\]

\begin{theorem}\label{thm:main}
The function \(f\) is finite continuous. Both raw suprema
\[
 p_{\C}(f)(v)=\sup\{g(v):g\in\C_R,\ g\leq f\},\qquad
 p_{\D}(f)(v)=\sup\{g(v):g\in\D_R,\ g\leq f\}
\]
equal the function
\[
 p(v)=\begin{cases}-a,&v(x)<+\infty,\\-b,&v(x)=+\infty.\end{cases}
\]
It is not upper semicontinuous. There is no greatest \(\C_R\)-minorant
and no greatest \(\D_R\)-minorant of \(f\).

For every integer \(m\geq1\), put \(A_m=\lceil ma\rceil\) and
\(B_m=\lceil mb\rceil\). The canonical ideal of the obstacle is exactly
\[
 \mathfrak a_m(f)=\{g\in R:v(g)+mf(v)\geq0\ \forall v\in X\}
       =\m^{B_m}+x\m^{A_m-1}.
\]
It is proper primary. Its actual geometric profile satisfies on every support
\[
 h_m=-\frac1m v(\mathfrak a_m(f))
    =\max\left\{-\frac{B_m}{m},-\frac{v(x)+A_m-1}{m}\right\},
 \qquad h_m\longrightarrow p\quad\text{pointwise},
\]
and has the exact error
\[
 \|p-h_m\|_\infty=\frac{\lceil mb\rceil}{m}-a\geq b-a.
\]
For \(a=1/2,b=1\), this error is \(1/2\) in every degree.
\end{theorem}

\section{All supports and the infinite-conductor fibres}

Since \(xy=0\), multiplicativity forces every semivaluation to kill
\(x\) or \(y\). Hence \(X=X_1\cup X_2\), a finite closed cover.
The first piece is the entire boundary of
\(T_1=k[x,\mathbf z]_{(x,\mathbf z)}\); the second is that of
\(T_2=k[y,\mathbf z]_{(y,\mathbf z)}\).
They are identified along the entire boundary of
\(B=k[\mathbf z]_{(\mathbf z)}\). These descriptions follow directly
by factoring a semivaluation through its killed coordinate, and preserve
every further prime support.

The normalization is \(S=T_1\times T_2\).
Indeed the natural embedding is a fibre product
\(R=T_1\times_B T_2\): pairs have equal restrictions at \(x=0,y=0\).
The polynomial fibre-product identity survives localization; common unit
denominators can be lifted from the two component polynomial rings.
The idempotent \((1,0)=x/(x+y)\) lies in the total quotient ring, since
\(x+y\) is a nonzerodivisor, and is integral over \(R\).
The component projections are surjective, so \(S=R+R(1,0)\) is finite.
Both component rings are normal, proving the asserted normalization.
Its conductor is
\[
 \mathfrak c=(x,y),\qquad R/\mathfrak c=B,
 \qquad S/\mathfrak c=B\times B.
\]
To check the conductor, multiply a pair by each component idempotent:
both resulting pairs descend precisely when their common restriction to
\(B\) is zero. This is the displayed ideal.

If \(v(x)\) is finite, only the first component can lift \(v\); if
\(v(y)\) is finite, only the second can. If both are infinite, there
are exactly two component lifts, with the same arbitrary semivaluation
on \(B\). Thus \(Z\), which is nonempty because \(d\geq2\), is
exactly the infinite-conductor stratum, and its fibres have two points.
No such support is deleted.

On \(X_2\), normalization gives
\(\min\{v(y),v(z_1),\ldots,v(z_{d-1})\}=1\), so \(v(y)\geq1\).
The reciprocal is continuous on the compact interval \([1,+\infty]\).
Consequently the second formula for \(f\) is continuous and takes values
in \([-b,-a]\); on \(Z\) it equals \(-a\).
The first formula is continuous and has the same value there.
The finite closed-cover pasting lemma proves continuity on \(X\).
This also proves full-fibre compatibility on the normalization.

\section{A geometric inequality on the second component}

Let \(o_i\) be the generic monomial valuation on \(T_i\) which assigns
value one to every regular coordinate. Its pullback to \(R\) kills
the other component coordinate, and belongs to \(X_i\).
For any normalized centred semivaluation \(w\) on \(T_i\), every
regular coordinate has value at least one. Therefore
\[
 w(g)\geq o_i(g)\quad\text{for every }g\in T_i.
\]
For completeness, represent \(g=P/Q\), where \(Q\) has nonzero
constant coefficient and hence value zero. Every monomial of \(P\)
has value at least its total degree, including monomials containing
killed coordinates. The ultrametric inequality proves the claim.
Generic monomial evaluation of \(P\) is its least total degree, so
no residue cancellation is assumed away for \(w\).

For every unit-or-primary \(R\)-ideal \(I\) this gives
\(w(I)\geq o_2(I)\) on \(X_2\); primary values are finite.
Taking negative finite nonnegative real combinations gives
\(q(w)\leq q(o_2)\) for \(q\in Q_R\).
Uniform limits preserve this inequality, as do pointwise infima of
actual decreasing nets. By the entire geometric comparison input,
\[
 g(w)\leq g(o_2)\qquad(w\in X_2,\ g\in\C_R\text{ or }\D_R).
\]
Since \(f(o_2)=-b\), every geometric minorant of \(f\) is at most
\(-b\) on the \emph{entire} second component, including \(Z\).
This is a statement about the original independently geometric classes,
not about a substitute ideal-order class.

\section{Exact canonical ideals and their positivity}

For \(g\in R\), write its component images as \(g_1,g_2\).
The first obstacle is constant, so its defining inequalities are equivalent to
\(g_1\in\m_1^{A_m}\). Necessity follows by testing \(o_1\);
sufficiency follows from the order bound on every semivaluation.
On the second component, testing \(o_2\) forces
\(g_2\in\m_2^{B_m}\). Conversely this condition implies, on every support,
\[
 w(g_2)\geq B_m\geq mb\geq
     m\left(a+\frac{b-a}{w(y)}\right),
\]
including \(w(y)=+\infty\). These are exactly all the obstacle tests.
The two conditions agree on the common quotient only for pairs whose
common restriction has order at least \(B_m\) in \(B\).

We verify that their inverse image is \(\m^{B_m}+x\m^{A_m-1}\).
This ideal has the required orders on both components, so one inclusion
is immediate. For the other, lift \(g_2\in\m_2^{B_m}\) from
the localized \(k[y,\mathbf z]\) subring of \(R\); this lift lies
in \(\m^{B_m}\) and has first component \(g_2(0,\mathbf z)\).
Subtract it from \(g\). The remaining second component is zero;
its first component is divisible by \(x\) and still belongs to
\(\m_1^{A_m}\), because \(B_m\geq A_m\).
Dividing by \(x\) in the regular polynomial local ring gives an
element of \(\m_1^{A_m-1}\), which lifts to
\(\m^{A_m-1}\) in \(R\). Multiplication by \(x\) kills its
second component. This proves the identity, with unit denominators
handled by the same localized subrings.

The ideal contains \(\m^{B_m}\) and is contained in \(\m\),
so it is proper primary. Ideal product and sum evaluations give the
formula for \(h_m\) in Theorem~\ref{thm:main}, with infinity retained.
Positivity is established before any limit is identified: on the actual
blowup \(Y_m=\operatorname{Bl}_{\mathfrak a_m(f)}\Spec R\), the pulled-back
ideal is \(\mathcal O_{Y_m}(-E_m)\), for an actual effective Cartier
divisor supported on the closed fibre, and this line bundle is relatively
ample. The Rees algebra is reduced as a subalgebra of \(R[t]\);
each generator chart embeds in its localization, where its ideal
generator is a nonzerodivisor. This proves actual Cartier effectivity
even though \(R\) is nonnormal and reducible.
The rational divisor \(-E_m/m\) is negative effective and its bundle
has positive degree on every contracted proper integral curve, by
ampleness restricted to that curve. It therefore belongs to the independently
geometric core and evaluates to \(h_m\). No arbitrary-model cofinality,
resolution, finite generation, or matrix positivity premise is used.

\section{Raw supremum, nonexistence, and exact uniform loss}

The constant \(-b=-b\,v(\m)\) is itself geometric and lies below \(f\).
The second-component inequality therefore makes either raw supremum
equal \(-b\) on \(X_2\).
On the remaining points \(v(x)<+\infty\), the obstacle is \(-a\)
and \(h_m(v)\to-a\). Since all \(h_m\) are geometric minorants,
both suprema equal \(-a\) there. This proves the asserted formula
and pointwise convergence on every support.

For each fixed degree the profiles lie between \(-B_m/m\) and \(p\).
On \(X_2\) the gap is \(B_m/m-b\), and on its complement it is
at most \(B_m/m-a\). A generic monomial point of the first component
with \(v(x)=r\geq B_m-A_m+1\) and \(v(z_j)=1\) attains the latter
gap. Thus the exact sup norm is \(B_m/m-a\), as claimed.

Let \(v_r\) be these first-component monomial points with integer
\(r\to+\infty\). They converge to \(v_\infty\), which kills
\(x,y\) and has generic monomial values one on all the \(z_j\).
Indeed for a polynomial \(F(x,\mathbf z)=\sum_i x^iF_i(\mathbf z)\),
\[
 v_r(F)=\min_i\{ri+\ord_{(\mathbf z)}F_i\}.
\]
If \(F_0\neq0\), this eventually equals its order; otherwise it tends
to infinity. Unit denominators have value zero. This proves convergence
on every element of \(R\), which is the boundary topology, and gives
\[
 \limsup_r p(v_r)=-a>-b=p(v_\infty).
\]
Hence \(p\) is not upper semicontinuous.

Any greatest \(\C_R\)- or \(\D_R\)-minorant would dominate
the constant \(-b\) and every \(h_m\), so the preceding pointwise
calculation would force it to equal \(p\). Both classes consist of
upper semicontinuous functions. This contradiction proves nonexistence.

Even ordinary upper semicontinuous regularization does not repair positivity.
Define \(p^*(v)=\inf_{U\ni v}\sup_U p\), over open neighbourhoods.
It is bounded upper semicontinuous, and \(p\leq p^*\leq f\) because
\(f\) is continuous. The sequence gives \(p^*(v_\infty)=-a\).
The open neighbourhood \(\{v:v(y)<2\}\) of \(o_2\) lies wholly in
\(X_2\), where \(p=-b\), so \(p^*(o_2)=-b\).
These values violate the mandatory geometric inequality
\(g(v_\infty)\leq g(o_2)\). Thus \(p^*\notin\C_R\cup\D_R\).

\section{Normalization and arbitrary primary gauges}

On the normalization, the first obstacle has envelope \(-a\), its
own constant geometric minorant. On the second component the same order
comparison makes its envelope the constant \(-b\).
Consequently the source envelope is the continuous geometric tuple
\((-a,-b)\), and it is not constant on any infinite-conductor fibre.
This proves failure of envelope preservation by normalization for this
continuous full-fibre-compatible obstacle. Exact class-level positive
descent in G331 is consistent: it requires the function being descended,
not merely its obstacle, to be constant on every fibre.
For this particular obstacle both normal components have continuous
geometric envelopes by the preceding elementary proof; no general normal
continuous-envelope theorem is needed.

For another proper primary gauge \(\mathfrak b\), set
\(d(v)=v(\mathfrak b)>0\) on \(X\) and \(T(v)=v/d(v)\).
The G331 homogeneous cocycle is an order bijection of both geometric
classes, with bounded positive continuous denominator. Transport the
obstacle by \(f_{\mathfrak b}(T(v))=f(v)/d(v)\).
Every canonical defining inequality is divided by the same \(d(v)\),
so the ideals are unchanged; profiles and the raw supremum divide by it.
The sequence \(T(v_r)\to T(v_\infty)\) retains the strict gap
\((b-a)/d(v_\infty)>0\). No greatest minorant exists in that gauge
either. Its uniform canonical error is bounded below by
\((b-a)/\max_X d\), by evaluating the same degree-dependent points.
This proves gauge-invariant failure without discarding supports.

\section{Source and scope record}

The sole incoming theorem is ENTIRE qualified G331, physically reread
here in bounded untruncated ranges 1--140, 141--270, 271--400, 401--530,
and 531--607, including all source and exact-scope annotations.
Its geometric comparison (n3 and n7--n8) is applied at this excellent
reduced polynomial local germ, with the actual primary-model and all-curve
conditions retained. Its primary-gauge cocycle is used only in the last
paragraph. This is fresh reading of the qualified grant, not fresh
recursive reading of its published primary sources.
Polynomial local rings being regular, normal and excellent, and their
ordinary monomial orders, are conventional algebraic background;
all normalization, conductor, support, obstacle, filtration and limit
computations needed here are proved above.

The proof is valid over every field, including imperfect and positive
characteristic fields; no integer coefficient is divided in the field.
It applies in every dimension at least two. It gives a nonnormal reduced
example, not a normal counterexample. In dimension one this conductor
quotient is zero dimensional and the centred quotient boundary is empty;
the construction then has no point \(v_\infty\), consistently with
geometric reduced-curve clipping. The boundary cone remains distinct
from literal classical or CLD logarithmic radial extensions.

General normal excellent higher-dimensional envelope regularity and
the full appropriate-existence classification remain open.
P12.4 products, P12.5 energy and P12.6 full operators remain separate.
Historical recursive/whole-book audits retain all G331 qualifications.
This complete scoped working proof is UNREVIEWED, with no original
closure or new accepted-paper claim. It was developed without inspection
of active native291 or native292 mathematics or hot feedback.
Compilation is separately unverified until an actual successful native
editor compiler result. The protected P16.4 editor remains open unchanged.
\end{document}
